%% kapa: exam=1 round=20 year=2009 period=2 form=A subject=appraisal_law subject_ko=부동산관계법규 seq=23 number=63
%% source: Q-Net 감정평가사 기출문제 (cst00309) artlSeq=5024178 file=extracted/1008/5024178/2009년 제20회 감정평가사 1차시험 기출문제/2교시.hwp
%% conversion: kordoc 4.24.51 -> tools/gen1.py (auto); review=proofread-sonnet; flags=none
\begin{question}{63}
부동산 가격공시 및 감정평가에 관한 법령상 감정평가사 또는 감정평가법인에 관한 설명으로 옳지 \uline{않은} 것은?
\begin{choices}
  \item 감정평가사는 감정평가업을 영위하기 위하여 2개 이상의 사무소를 설치할 수 있다.
  \item 감정평가사는 그 업무를 효율적으로 수행하고 공신력을 높이기 위하여 3인 이상의 감정평가사로 구성된 합동사무소를 설치할 수 있다.
  \item 감정평가법인에는 10인 이상의 감정평가사를 두어야 한다.
  \item 감정평가법인은 사원 전원의 동의 또는 주주총회의 의결이 있는 때에는 국토해양부장관의 인가를 받아 다른 감정평가법인과 합병할 수 있다.
  \item 감정평가법인은 당해 법인의 소속감정평가사 외의 자로 하여금 감정평가업자의 업무를 하게 하여서는 아니 된다.
\end{choices}
\end{question}
\answer{1}
